%==============================================================================
\section{The \Onehop Operator}
\label{sec:onehop}
%==============================================================================

As established in \Cref{sec:background:graph}, \system maintains two copies of
each edge label: a \emph{forward list} sorted by source node id and a
\emph{reverse list} sorted by destination node id. The \onehop operator exploits
both copies to evaluate the pattern $n_i \rightarrow e \rightarrow n_j$, joining a
source node label $S\,(id, [\mathit{prop}])$, an edge label $E$, and a
destination label $D$ on their foreign-key relationship. The result is equivalent
to the three-way equi-join $S \bowtie E \bowtie D$, and its size is exactly
$|E|$, a publicly known quantity in our threat model. Because the forward list
already groups edges by source and the reverse list groups them by destination,
the source-side and destination-side joins can be processed in parallel as two
independent passes.

An oblivious \onehop must hide three sources of leakage. First,
\emph{hash-bucket access}: the memory pattern of retrieving a node's attributes
by id would expose which key was probed. Second, \emph{node degree}: the number
of source-side (or destination-side) lookups for a given id would equal its
degree, leaking the degree distribution of the graph. Third, \emph{predicate
selectivity}: the number of edges that satisfy the query predicates would reveal
how many tuples match. The algorithm below neutralizes the first two by
construction; the third is handled by always emitting $|E|$ rows and deferring
any selectivity-dependent sizing, as we discuss at the end of the section.

%------------------------------------------------------------------------------
\subsection{Online Algorithm}
\label{sec:onehop:algorithm}
%------------------------------------------------------------------------------

Algorithm~\ref{alg:onehop} presents the operator. The two node indices $H_S$ and
$H_D$ are oblivious hash tables built offline from the node tables, keyed by node
id, once per node label; this query-independent setup is amortized across all
queries against the same graph. We use the oblivious hash table
of~\cite{h2o2ram}, whose single defining constraint shapes our design:
\emph{each distinct non-dummy key may be queried at most once} --- the table
marks an entry once probed, and a repeat query on the same key breaks its
obliviousness guarantee. Since many edges typically share a source (or
destination) node, a one-hop join over the edge table would otherwise issue many
queries for the same key. The phases that probe the index are structured
precisely to satisfy this single-query-per-key constraint while still attaching
node attributes to every edge.

\begin{algorithm}[t]
\caption{Oblivious \onehop Operator}
\label{alg:onehop}
\begin{algorithmic}[1]
\Require Forward list $E_{\!f}$ (sorted by $src\_id$), reverse list $E_{\!r}$
         (sorted by $dst\_id$); precomputed indices $H_S$, $H_D$; optional
         predicates $P$
\Ensure  Result $R$ with $|R| = |E|$
\Statex
\State \textbf{Source branch} \Comment{concurrent with the destination branch}
\State $\textsc{Dedup}(E_{\!f})$: for each row $i$,\algoIndent{1}
       $E_{\!f}[i].k \gets \textsc{OCmpSet}\big(E_{\!f}[i].src\_id = E_{\!f}[i\!-\!1].src\_id,\ \bot,\ E_{\!f}[i].src\_id\big)$
\State for each row $i$: probe $H_S[E_{\!f}[i].k]$ \Comment{all $|E|$ rows; dummy keys are no-ops}
\State $\textsc{Redup}(E_{\!f})$: for each row $i$,\algoIndent{1}
       $E_{\!f}[i].sattr \gets \textsc{OCmpSet}\big(E_{\!f}[i].k = \bot,\ E_{\!f}[i\!-\!1].sattr,\ E_{\!f}[i].sattr\big)$
\State $\textsc{Merge}(E_{\!f}, \text{source attributes})$ \Comment{attach source columns}
\Statex
\State \textbf{Destination branch} \Comment{symmetric over $dst\_id$, $H_D$}
\State $\textsc{Dedup}$ / probe $H_D$ / $\textsc{Redup}$ / $\textsc{Merge}$ on $E_{\!r}$
\State $E_{\!r} \gets \textsc{OSort}(E_{\!r})$ into edge order \Comment{re-align to source branch}
\Statex
\State \textbf{Combine}
\State $R \gets \textsc{Merge}(E_{\!f}, E_{\!r})$ \Comment{column-wise, aligned per edge}
\If{$P \neq \emptyset$} \Comment{standalone query only; skipped under decomposition}
    \State $R \gets \textsc{OCompact}\big(\text{rows of } R \text{ satisfying } P\big)$
\EndIf
\State \Return $R$
\end{algorithmic}
\end{algorithm}

\paragraph{Two branches over the pre-sorted lists.}
Because the forward and reverse edge lists are already maintained in sorted order
(\Cref{sec:background:graph}), the operator needs no online sort to group edges by
key --- the grouping comes for free from the storage layout. The source branch
consumes the forward list and the destination branch the reverse list, and the
two run concurrently on independent working copies.

\paragraph{Deduplicate, probe, reduplicate.}
Within a branch, \emph{deduplication} enforces the single-query-per-key
constraint: in a linear scan over the sorted list, we compare each row's key with
its predecessor and, using \textsc{OCmpSet}, obliviously replace it with a dummy
($\bot$) whenever the two match. Exactly one row per distinct key --- the first in
its run --- keeps a live key. We then probe the index \emph{once for every row},
dummies included; a dummy key resolves to a no-op the table absorbs, while live
keys retrieve their node attributes. Probing every row, rather than only the live
ones, keeps the query count equal to $|E|$ regardless of the graph's degree
distribution; probing only live keys would let the query count leak the number of
distinct sources. \emph{Reduplication} then refills the suppressed rows: a second
scan copies, via \textsc{OCmpSet}, the probed attributes from the preceding row
whenever the current key is a dummy, propagating each run's attributes forward. A
\textsc{Merge} attaches the retrieved node attributes to the edge rows as new
columns.

\paragraph{Align and combine.}
The source branch finishes in $src\_id$ order, but the destination branch finishes
in $dst\_id$ order; we therefore oblivious-sort the destination branch back into
the same per-edge order so the two $|E|$-row tables align row-for-row. A final
\textsc{Merge} then combines them column-wise, giving every edge both its source
and destination attributes. This single re-alignment is the operator's only online
sort.

\paragraph{Filtering.}
When the operator answers a standalone one-hop query, we evaluate the predicates
$P$ over every row after the combine and obliviously compact the matches; the
result size is the query's (public) output size. When the operator instead serves
as a building block under query decomposition (\Cref{sec:decomposition}), no
predicates are applied here --- the output remains exactly $|E|$ rows, keeping
intermediate sizes independent of selectivity.

%------------------------------------------------------------------------------
% Alternatives considered (closes the section).
%------------------------------------------------------------------------------
The one-hop pattern $S \bowtie E \bowtie D$ could in principle be evaluated as two
sequential oblivious binary joins using an existing operator such as
Obliviator~\cite{mavrogiannakis2025obliviator}. Such a composition is
security-equivalent --- it reveals only the same public sizes --- but it pays for
two full oblivious-join executions, each performing its own oblivious sorts over
the joined relations, and it materializes an intermediate result that the second
join must then re-sort. Our operator instead enriches the edge table in place from
both directions, replacing that machinery with one pass of probes per side and a
single re-alignment sort. \Cref{sec:evaluation} shows this yields a
$1.1$--$1.7\times$ latency improvement over the two-join baseline at equal
security.

%------------------------------------------------------------------------------
% \subsection{Analysis}
% \label{sec:onehop:analysis}
% Correctness + obliviousness arguments. Deferred until \S\S4--6 are stable.
%------------------------------------------------------------------------------
